\chapter{Resultados}
\label{cap:resultados}

\instruccion{Reporte toda la evidencia obtenida, incluidos resultados nulos,
negativos o inesperados. La organización debe seguir la estrategia de
evaluación definida en metodología y facilitar el contraste de la hipótesis y
el cumplimiento de los objetivos. No es obligatorio crear una subsección por
pregunta de investigación ni presentar el capítulo como una bitácora del
desarrollo. Separe la evidencia de su interpretación amplia, que corresponde al
capítulo de discusión.}

\section{Evidencia obtenida}

Describa brevemente qué evaluación se ejecutó y qué evidencia quedó disponible:
participantes, casos, conjuntos de datos, ejecuciones, versiones, escenarios,
instrumentos o registros. Reporte exclusiones y sus causas, datos faltantes y
características necesarias para interpretar los resultados. No revele
información que permita identificar personas sin autorización.

\section{Resultados de la evaluación}

Presente tablas, figuras, estimaciones, intervalos, distribuciones, categorías,
citas de participantes o evidencia técnica según el diseño. Organice esta
sección por dimensión evaluada, iteración, experimento, método comparado,
escenario o criterio de aceptación: elija la estructura que haga visible la
relación entre la evidencia y lo prometido en la metodología. Toda figura o
tabla debe citarse y explicarse en el texto.

\instruccion{Seleccione la ruta de evaluación apropiada. En software dirigido a
personas, incluya validación con usuarios representativos mediante tareas y un
instrumento o escala justificados; reporte tamaño y características de la
muestra, condiciones de aplicación, puntuaciones, dispersión y comentarios
relevantes. En comparaciones de métodos, algoritmos, CRISP-DM o KDD, reporte
conjuntos de datos, particiones, líneas base, métricas, repeticiones,
incertidumbre y análisis de errores. En evaluaciones técnicas, use criterios
medibles o escenarios de atributos de calidad para rendimiento, seguridad,
confiabilidad, mantenibilidad u otras propiedades. Combine rutas cuando el
artefacto tenga tanto interacción humana como comportamiento técnico.}

\ifmostrarEjemplos
La \cref{tab:ejemplo-validacion-usuarios} ilustra una forma compacta de resumir
una validación con usuarios. Los nombres y valores son ficticios: reemplácelos
por el instrumento, la escala y los datos reales del estudio.

\begin{table}[H]
\centering
\caption{Ejemplo de resultados de validación con usuarios}
\label{tab:ejemplo-validacion-usuarios}
\small
\begin{tabularx}{\textwidth}{L{3.1cm}L{2.5cm}L{2.3cm}Y}
\toprule
\textbf{Dimensión} & \textbf{Instrumento o escala} &
\textbf{Resultado} & \textbf{Criterio de interpretación} \\
\midrule
Usabilidad percibida & Escala validada seleccionada &
Media y dispersión & Umbral o guía definida antes del análisis \\
Satisfacción & Ítems y rango de respuesta &
Distribución de respuestas & Criterio establecido en metodología \\
Eficacia en tareas & Registro de tareas &
Porcentaje de éxito & Meta de aceptación del sistema \\
\bottomrule
\end{tabularx}
\end{table}
\fi

\section{Contraste de la hipótesis y cumplimiento de objetivos}

Integre la evidencia necesaria para establecer si la hipótesis recibió apoyo y
en qué grado se alcanzó cada objetivo. Si la hipótesis es estadística, reporte
la prueba o modelo, supuestos, estimación, incertidumbre, tamaño del efecto y la
decisión sobre la hipótesis nula. Si se formuló una hipótesis de diseño o una
proposición, confróntela con criterios de aceptación y fuentes de evidencia
previamente definidos. Evite afirmar que una hipótesis fue ``demostrada'' y no
confunda ausencia de evidencia con evidencia de ausencia.

\instruccion{Las preguntas de investigación siguen siendo una referencia para
comprobar que no quedó ninguna sin atender, pero no tienen que convertirse en
subtítulos. Los hallazgos adicionales pueden incorporarse al final de la sección
correspondiente, identificándolos como exploratorios cuando no estaban previstos
antes de observar los datos.}
